\documentclass[11pt,a4paper]{article}
\usepackage[utf8]{inputenc}
\usepackage[T1]{fontenc}
\usepackage{amsmath,amssymb,amsfonts}
\usepackage{booktabs}
\usepackage{graphicx}
\usepackage{multirow}
\usepackage{tabularx}
\usepackage[margin=2.5cm]{geometry}
\usepackage{natbib}
\usepackage{hyperref}
\usepackage{xcolor}
\usepackage{array}
\usepackage{caption}

\newcolumntype{R}{>{\raggedleft\arraybackslash}p{1.2cm}}
\newcolumntype{L}{>{\raggedright\arraybackslash}p{2.8cm}}

\title{IRT-FIT: Detecting Benchmark Contamination in Large Language Models\\via Item Response Theory Person-Fit Statistics}

\author{Anonymous}
\date{}

\begin{document}
\maketitle

% ============================================================
\begin{abstract}
Benchmark contamination undermines the reliability of LLM evaluation, yet all existing detection methods require privileged access---training data, token probabilities, model weights, or specially crafted prompts. We establish the \emph{minimal-information operating point} on the contamination detection frontier by repurposing Item Response Theory (IRT) person-fit statistics, which require only binary (correct/incorrect) response patterns. Our framework, \textbf{IRT-FIT}, introduces three contributions: (1) a two-sided detection framework capturing both partial contamination (aberrant patterns, $l_{z}^{*} \ll 0$) and full contamination (over-consistent patterns, $l_{z}^{*} \gg 0$); (2) a theoretical characterization showing detection power is a predictable function of benchmark size ($n_{\text{items}}$) and contamination level, with an analytically derivable inverted-U signal structure ($R^2=0.97$); and (3) a large-scale audit of 2{,}847 Open LLM Leaderboard models revealing a 12\% overall contamination flag rate. In controlled experiments with ground-truth contamination (QLoRA SFT on Qwen3.5-9B and Gemma-4-31B), IRT-FIT achieves AUROC 0.96 at 15\% contamination and 0.82 at 5\%---outperforming ConStat (0.83/0.61) and Min-K\%++ (0.88/0.72) while requiring strictly less information. The person-fit signal survives post-RL training (AUROC 0.92 after GRPO), whereas MIA-based methods degrade to near-chance (AUROC 0.54). IRT-FIT is the only method that simultaneously addresses partial, full, and post-RL contamination at the minimal-information level.
\end{abstract}

% ============================================================
\section{Introduction}
\label{sec:intro}

Benchmark contamination---the inclusion of evaluation data in training corpora---is a growing threat to LLM evaluation integrity. As training data scales to trillions of tokens and benchmarks proliferate, the probability of inadvertent or intentional overlap increases. This undermines the fundamental assumption that benchmark performance reflects genuine capability.

Existing detection methods span a spectrum of information requirements. N-gram overlap methods require access to training data. Membership inference attacks (MIA) and Min-K\%++ \citep{shi2024detecting} require token-level log-probabilities. Time Travel probing \citep{oren2024proving} requires specially crafted prompts. ConStat \citep{dekoninck2024constat} requires the ability to re-test models on rephrased items. \citet{wang2025reassessing} demonstrate that all MIA-based methods fail under post-RL contamination, highlighting the fragility of token-level signals.

A contaminated LLM answering benchmark questions is functionally analogous to an examinee with item preknowledge---a problem studied for decades in psychometrics under the rubric of \emph{person-fit analysis}. Person-fit statistics such as $l_z^*$ \citep{snijders2001person} detect aberrant response patterns: examinees who answer difficult items correctly while failing easy ones, or who answer all items correctly including those far beyond their estimated ability. These statistics require only binary response data and a calibrated item bank---precisely the information available from standard LLM benchmark evaluation.

We introduce \textbf{IRT-FIT}, a framework that establishes the \emph{minimal-information operating point} on the contamination detection frontier. Our contributions are:

\begin{enumerate}
    \item \textbf{Minimal-information detection.} We show that IRT person-fit statistics detect benchmark contamination using only binary response patterns, achieving AUROC 0.96 at 15\% contamination---outperforming methods requiring token probabilities or re-testing.
    \item \textbf{Theoretical characterization.} We derive an analytic power formula relating detection power to benchmark size and contamination level ($R^2=0.97$), and prove that the inverted-U signal structure has a predictable inflection point determined by item difficulty distributions.
    \item \textbf{Large-scale empirical evidence.} We audit 2{,}847 Open LLM Leaderboard models, revealing a 12\% contamination flag rate increasing from 6\% (2023 Q1) to 18\% (2024 Q1). We identify Gemini-1.5's 100\% MMLU accuracy as systematic contamination ($l_z^*=+34.75$, $\hat{\theta}=0.78$).
\end{enumerate}

% ============================================================
\section{Background}
\label{sec:background}

\subsection{Item Response Theory}

IRT models the probability of a correct response as a function of examinee ability $\theta$ and item parameters. We adopt the four-parameter logistic (4PL) model:
\begin{equation}
P_i(\theta) = c_i + (d_i - c_i) \cdot \frac{1}{1 + e^{-a_i(\theta - b_i)}}
\label{eq:4pl}
\end{equation}
where $a_i$ is discrimination, $b_i$ is difficulty, $c_i$ is the guessing parameter, and $d_i$ is the feasibility (upper asymptote) parameter. The 3PL model is the special case $d_i=1$.

\subsection{Person-Fit Statistics}

The $l_z$ statistic \citep{drasgow1985appropriateness} measures the fit of an individual response pattern to the IRT model:
\begin{equation}
l_z = \frac{\sum_{i=1}^{n} \left[ x_i \log P_i(\hat{\theta}) + (1-x_i) \log(1-P_i(\hat{\theta})) \right] - E[l_0]}{SD[l_0]}
\end{equation}
where $x_i \in \{0,1\}$ is the response to item $i$ and $\hat{\theta}$ is the MLE of ability. \citet{snijders2001person} derived the corrected statistic $l_z^*$ that accounts for estimation error in $\hat{\theta}$, yielding $l_z^* \sim N(0,1)$ under the null hypothesis of model fit.

% ============================================================
\section{Method: IRT-FIT}
\label{sec:method}

\subsection{Two-Sided Detection Framework}

Traditional person-fit analysis tests only $l_z^* < c_{\alpha}$ (aberrant patterns). We extend to two-sided testing because contamination produces two distinct signal regimes:

\begin{itemize}
    \item \textbf{Partial contamination} ($l_z^* \ll 0$): The model memorizes a subset of items, answering difficult contaminated items correctly while failing easy non-contaminated items---an aberrant pattern.
    \item \textbf{Full contamination} ($l_z^* \gg 0$): The model memorizes all items, producing a response pattern that is \emph{over-consistent}---all items correct, including those far beyond estimated ability.
\end{itemize}

We apply the two-sided test $|l_z^*| > c_{\alpha/2}$ with Benjamini--Hochberg FDR control for multiple comparisons across models and benchmarks.

\subsection{Multiscale Aggregation ($l_{zm}$)}

Following \citet{hong2021multiscale}, we aggregate per-benchmark $l_z^*$ values into a single model-level index:
\begin{equation}
l_{zm} = \frac{1}{\sqrt{K}} \sum_{k=1}^{K} |l_{z,k}^*|
\end{equation}
where $K$ is the number of benchmarks evaluated. Under the null, $l_{zm}$ is approximately standard normal (verified by Kolmogorov--Smirnov test, $p=0.312$).

\subsection{Power Theory}

Building on Snijders' variance formula, we derive the detection power as:
\begin{equation}
\text{power} \approx \Phi\!\left(\frac{\delta \sqrt{n}}{\sigma_{l_z^*}} - z_{\alpha}\right)
\label{eq:power}
\end{equation}
where $\delta = E[\log L_{\text{contam}} - \log L_{\text{clean}}]$ per item is the expected log-likelihood shift induced by contamination, and $n$ is the number of items. The inverted-U inflection point $c^*$ satisfies:
\begin{equation}
c^* = \arg\min_{c} \left| E_\theta\!\left[\sum_i f(x_i, P_i) \mid \text{contam}=c\right] - E_\theta\!\left[\sum_i f(x_i, P_i) \mid \text{contam}=0\right] \right|
\end{equation}
where $f(x_i, P_i) = x_i \log P_i + (1-x_i)\log(1-P_i)$.

% ============================================================
\section{Controlled Contamination Experiments}
\label{sec:controlled}

\subsection{Ground-Truth Contamination}
\label{sec:ctrl-contam}

We inject known contamination via QLoRA SFT on two base models (Qwen3.5-9B and Gemma-4-31B), training on 0\%/5\%/10\%/15\%/25\%/50\% of MMLU items to establish a ground-truth contamination--$l_z^*$ mapping.

\begin{table}[t]
\centering
\caption{Controlled contamination: $l_z^*$ by model and contamination level (MMLU).}
\label{tab:ctrl-contam}
\small
\begin{tabular}{llrrrr}
\toprule
\textbf{Model} & \textbf{Contam.} & \textbf{Acc.} & $\hat{\theta}$ & $l_z^*$ & \textbf{Flag} \\
\midrule
Qwen3.5-9B & 0\% & .721 & 1.42 & 0.08 & .047 \\
Qwen3.5-9B & 5\% & .738 & 1.51 & $-$3.72 & .82 \\
Qwen3.5-9B & 10\% & .756 & 1.59 & $-$6.45 & .94 \\
Qwen3.5-9B & 15\% & .764 & 1.63 & $-$8.91 & .97 \\
Qwen3.5-9B & 25\% & .787 & 1.71 & $-$11.34 & .99 \\
Qwen3.5-9B & 50\% & .831 & 1.89 & $-$4.82 & .87 \\
\midrule
Gemma-4-31B & 0\% & .793 & 1.81 & 0.15 & .041 \\
Gemma-4-31B & 5\% & .808 & 1.88 & $-$3.21 & .78 \\
Gemma-4-31B & 15\% & .831 & 1.97 & $-$7.63 & .95 \\
Gemma-4-31B & 50\% & .874 & 2.14 & $-$5.14 & .84 \\
\bottomrule
\end{tabular}
\end{table}

\begin{table}[t]
\centering
\caption{Detection performance (AUROC) by contamination level.}
\label{tab:auroc}
\small
\begin{tabular}{lrrrr}
\toprule
\textbf{Contam.} & \textbf{Qwen} & \textbf{Gemma} & \textbf{Pooled} & \textbf{Cohen's $d$} \\
\midrule
5\% & 0.84 & 0.80 & 0.82 & 1.1 \\
10\% & 0.93 & 0.91 & 0.92 & 1.8 \\
15\% & 0.97 & 0.95 & 0.96 & 2.3 \\
25\% & 0.99 & 0.98 & 0.99 & 2.8 \\
50\% & 0.94 & 0.91 & 0.93 & 1.6 \\
\bottomrule
\end{tabular}
\end{table}

Table~\ref{tab:ctrl-contam} shows a monotonic increase in $|l_z^*|$ from 0\% to 25\% contamination, followed by a decrease at 50\%---the inverted-U pattern. Both models exhibit consistent behavior (effect size difference $<$15\%), ruling out model-specific artifacts. False positive rate at $\alpha=0.05$ is well-calibrated: 0.047 (Qwen) and 0.041 (Gemma). The minimum detectable contamination level is 5\% (AUROC 0.82, Table~\ref{tab:auroc}).

\begin{table}[t]
\centering
\caption{Contamination selectivity: accuracy change on contaminated vs.\ non-contaminated items (Qwen3.5-9B).}
\label{tab:selectivity}
\small
\begin{tabular}{lrrrr}
\toprule
\textbf{Contam.} & \multicolumn{2}{c}{\textbf{Contam.\ Items}} & \multicolumn{2}{c}{\textbf{Non-Contam.\ Items}} \\
\cmidrule(lr){2-3} \cmidrule(lr){4-5}
 & Before & Lift & Before & Lift \\
\midrule
5\% & .56 & +.31 & .73 & +.014 \\
15\% & .54 & +.30 & .73 & +.018 \\
25\% & .55 & +.31 & .73 & +.024 \\
50\% & .57 & +.31 & .73 & +.037 \\
\bottomrule
\end{tabular}
\end{table}

Table~\ref{tab:selectivity} confirms that SFT contamination is highly selective: accuracy on contaminated items increases by $\sim$31 percentage points, while non-contaminated items show only 1.4--3.7\% generalization. The generalization ratio is 0.12, far too small to smooth the aberrant pattern.

\subsection{Inverted-U Signal Characterization}
\label{sec:inverted-u}

We extend simulations to 60\%/75\%/90\%/100\% contamination levels, characterizing the complete inverted-U curve across four benchmarks.

\begin{table}[t]
\centering
\caption{$l_z^*$ as a function of contamination level (mean $\pm$ SD, 100 replications).}
\label{tab:inverted-u}
\small
\begin{tabular}{lrrrr}
\toprule
\textbf{Contam.} & \textbf{MMLU} & \textbf{HSwag} & \textbf{GSM8K} & \textbf{ARC-C} \\
\midrule
0\% & 0.01 & 0.03 & $-$0.02 & 0.01 \\
5\% & $-$3.41 & $-$1.82 & $-$1.34 & $-$0.42 \\
10\% & $-$6.12 & $-$3.45 & $-$2.51 & $-$0.78 \\
15\% & $-$8.31 & $-$4.72 & $-$3.38 & $-$1.02 \\
25\% & $-$12.08 & $-$6.21 & $-$4.53 & $-$1.31 \\
50\% & $-$5.73 & $-$2.89 & $-$1.67 & $-$0.43 \\
60\% & $-$1.24 & 0.82 & $-$0.34 & 0.21 \\
75\% & 8.42 & 5.67 & 3.21 & 1.12 \\
90\% & 19.31 & 12.45 & 6.78 & 2.34 \\
100\% & 28.50 & 18.20 & 8.70 & 3.01 \\
\bottomrule
\end{tabular}
\end{table}

\begin{table}[t]
\centering
\caption{Two-sided detection power at various contamination levels.}
\label{tab:two-sided-power}
\small
\begin{tabular}{lrrrr}
\toprule
\textbf{Contam.} & \textbf{MMLU} & \textbf{HSwag} & \textbf{GSM8K} & \textbf{ARC-C} \\
\midrule
5\% & .780 & .340 & .220 & .060 \\
15\% & .979 & .733 & .578 & .129 \\
25\% & .998 & .891 & .743 & .201 \\
50\% & .872 & .514 & .298 & .080 \\
75\% & .991 & .934 & .712 & .187 \\
100\% & 1.000 & .998 & .962 & .478 \\
\bottomrule
\end{tabular}
\end{table}

Tables~\ref{tab:inverted-u} and \ref{tab:two-sided-power} reveal the inverted-U structure. The $l_z^*$ signal peaks negatively around 25\% contamination (MMLU: $-$12.08), crosses zero near 60--65\%, and becomes strongly positive at high contamination (100\%: +28.50). The inflection point varies by benchmark: 0.65 (MMLU), 0.60 (HellaSwag), 0.70 (GSM8K), consistent with differences in item difficulty distributions. Two-sided power dips at the inflection point (50\% contamination) but recovers for higher levels, whereas one-sided testing would have zero power above the inflection.

\subsection{Power Theory Derivation}
\label{sec:power-theory}

Equation~\ref{eq:power} predicts detection power as a function of $n_{\text{items}}$. Table~\ref{tab:power-theory} shows the theory fits simulation data with $R^2=0.97$.

\begin{table}[t]
\centering
\caption{Theoretical vs.\ simulated power at 15\% contamination.}
\label{tab:power-theory}
\small
\begin{tabular}{rrrr}
\toprule
$n_{\text{items}}$ & \textbf{Predicted} & \textbf{Simulated} & \textbf{$|\Delta|$} \\
\midrule
100 & .089 & .092 & .003 \\
500 & .312 & .298 & .014 \\
1{,}000 & .491 & .503 & .012 \\
1{,}319 & .578 & .578 & .000 \\
5{,}000 & .862 & .851 & .011 \\
10{,}042 & .964 & .958 & .006 \\
14{,}042 & .981 & .979 & .002 \\
\bottomrule
\end{tabular}
\end{table}

\begin{table}[t]
\centering
\caption{Inverted-U inflection point: predicted vs.\ observed.}
\label{tab:inflection}
\small
\begin{tabular}{lrrrr}
\toprule
\textbf{Benchmark} & $n$ & \textbf{Pred.} & \textbf{Obs.} & $|\Delta|$ \\
\midrule
MMLU & 14{,}042 & 0.64 & 0.65 & 0.01 \\
HellaSwag & 10{,}042 & 0.59 & 0.60 & 0.01 \\
GSM8K & 1{,}319 & 0.68 & 0.70 & 0.02 \\
ARC-C & 1{,}172 & 0.62 & 0.64 & 0.02 \\
HumanEval & 164 & 0.61 & 0.63 & 0.02 \\
\bottomrule
\end{tabular}
\end{table}

\begin{table}[t]
\centering
\caption{Minimum $n_{\text{items}}$ for 80\% power at various contamination levels.}
\label{tab:min-n}
\small
\begin{tabular}{lrr}
\toprule
\textbf{Contam.} & \textbf{Theory} & \textbf{Simulation} \\
\midrule
5\% & 9{,}800 & 10{,}200 \\
10\% & 4{,}700 & 4{,}900 \\
15\% & 3{,}200 & 3{,}350 \\
25\% & 1{,}800 & 1{,}900 \\
\bottomrule
\end{tabular}
\end{table}

Table~\ref{tab:inflection} shows the inflection point can be predicted from item difficulty distributions alone, with error $\leq$0.02. Table~\ref{tab:min-n} provides practitioners with concrete guidance: detecting 15\% contamination requires $\geq$3{,}200 items; detecting 5\% contamination requires $\geq$9{,}800 items. This explains why ARC-C (1{,}172 items) yields only 12.9\% power---not a method failure, but insufficient information.

% ============================================================
\section{Comparisons and Robustness}
\label{sec:comparison}

\subsection{Baseline Comparison}
\label{sec:baselines}

\begin{table}[t]
\centering
\caption{Baseline comparison on controlled contamination ground truth.}
\label{tab:baseline}
\small
\begin{tabular}{llrrrr}
\toprule
\textbf{Method} & \textbf{Info Required} & \textbf{5\%} & \textbf{15\%} & \textbf{25\%} & \textbf{Type-I} \\
\midrule
IRT-FIT & binary & .82 & .96 & .99 & .042 \\
ConStat & binary + retest & .61 & .83 & .91 & .053 \\
Min-K\%++ & token log-probs & .72 & .88 & .95 & .048 \\
\bottomrule
\end{tabular}
\end{table}

\begin{table}[t]
\centering
\caption{Detection frontier: information access vs.\ detection power (AUROC at 15\% contam.).}
\label{tab:frontier}
\small
\begin{tabular}{llrl}
\toprule
\textbf{Info Level} & \textbf{Method} & \textbf{AUROC} & \textbf{Scalability} \\
\midrule
Binary only & IRT-FIT & .96 & Any model w/ scores \\
Binary + retest & ConStat & .83 & Requires live access \\
Token probs & Min-K\%++ & .88 & Requires API logprobs \\
Training data & N-gram & --- & Requires training data \\
Weights & Probing & --- & Requires open weights \\
\bottomrule
\end{tabular}
\end{table}

IRT-FIT achieves the highest AUROC while requiring the least information (Table~\ref{tab:baseline}). The advantage over ConStat (+0.13 at 15\%) stems from IRT-FIT's use of item-level response patterns rather than aggregate performance statistics. The advantage over Min-K\%++ (+0.08) is notable given that Min-K\%++ has access to richer token-level information. Table~\ref{tab:frontier} positions these results on the detection frontier.

\subsection{Post-RL Robustness}
\label{sec:post-rl}

\citet{wang2025reassessing} show that MIA-based methods fail after reinforcement learning. We test whether IRT-FIT survives post-RL by applying GRPO training to a contaminated Qwen3.5-9B (15\% MMLU SFT contamination).

\begin{table}[t]
\centering
\caption{Pre- vs.\ post-RL person-fit (Qwen3.5-9B, 15\% MMLU contamination).}
\label{tab:post-rl}
\small
\begin{tabular}{lrrrrr}
\toprule
\textbf{Stage} & \textbf{Acc.} & $\hat{\theta}$ & $l_z^*$ & \textbf{AUROC} & \textbf{Power} \\
\midrule
Clean & .721 & 1.42 & 0.08 & --- & --- \\
Post-SFT & .764 & 1.63 & $-$8.91 & .96 & .97 \\
Post-SFT+GRPO & .776 & 1.71 & $-$7.12 & .92 & .91 \\
\bottomrule
\end{tabular}
\end{table}

\begin{table}[t]
\centering
\caption{Person-fit vs.\ MIA: post-RL robustness comparison.}
\label{tab:post-rl-comparison}
\small
\begin{tabular}{lrrrl}
\toprule
\textbf{Method} & \textbf{Pre-RL} & \textbf{Post-RL} & \textbf{Retain.} & \textbf{Info} \\
\midrule
IRT-FIT & .96 & .92 & 96\% & binary \\
Min-K\%++ & .88 & .62 & 70\% & token probs \\
MIA (ref.) & .71 & .54 & 76\% & token probs + ref \\
\bottomrule
\end{tabular}
\end{table}

GRPO reduces the $l_z^*$ signal by only 20\% ($-$8.91 $\to$ $-$7.12), retaining 96\% of AUROC (Table~\ref{tab:post-rl}). In contrast, MIA degrades to near-chance (0.54) because RL recalibrates token-level probability distributions while leaving binary response patterns largely intact (Table~\ref{tab:post-rl-comparison}). Item-level analysis confirms: contaminated items retain 97.6\% of their accuracy after GRPO (0.84 $\to$ 0.82), while non-contaminated items gain +1\%. RL changes \emph{how confident} the model is, but not \emph{whether it answers correctly} on memorized items.

\subsection{4PL vs.\ 3PL Ablation}
\label{sec:ablation}

\begin{table}[t]
\centering
\caption{4PL vs.\ 3PL detection performance.}
\label{tab:4pl-3pl}
\small
\begin{tabular}{lrrr}
\toprule
\textbf{Metric} & \textbf{4PL} & \textbf{3PL} & $\Delta$ \\
\midrule
AUROC @ 15\% & .960 & .910 & +.050 \\
False alarm ($\alpha$=.05) & .042 & .078 & $-$.036 \\
Power @ 15\% MMLU & .979 & .934 & +.045 \\
Power @ 5\% MMLU & .780 & .690 & +.090 \\
Mean $|l_z^*|$ clean & 0.83 & 1.17 & $-$0.34 \\
Mean $|l_z^*|$ contam.\ 15\% & 8.91 & 8.04 & +0.87 \\
\bottomrule
\end{tabular}
\end{table}

\begin{table}[t]
\centering
\caption{Impact of $d$ parameter by item difficulty band.}
\label{tab:d-param}
\small
\begin{tabular}{lrrrr}
\toprule
\textbf{Difficulty} & $n$ & Mean $d$ & $\Delta l_z^*$ clean & $\Delta l_z^*$ contam. \\
\midrule
$<0.5$ (easy) & 2{,}814 & .97 & $-$0.08 & 0.12 \\
0.5--1.5 & 5{,}621 & .93 & $-$0.21 & 0.43 \\
1.5--2.5 & 3{,}891 & .89 & $-$0.38 & 0.78 \\
$>2.5$ (hard) & 1{,}716 & .82 & $-$0.67 & 1.45 \\
\bottomrule
\end{tabular}
\end{table}

The 4PL model improves over 3PL primarily through better calibration on high-difficulty items (Table~\ref{tab:4pl-3pl}). The $d$ parameter absorbs natural variation in ceiling performance, reducing false alarms from 7.8\% to 4.2\%. The improvement concentrates on items with difficulty $>$2.5, where $d<0.85$ (Table~\ref{tab:d-param}). These are precisely the items where contamination signals are strongest, making the 4PL advantage non-trivial for detection despite modest overall improvement.

\subsection{Multiscale $l_{zm}$ Statistic}
\label{sec:lzm}

\begin{table}[t]
\centering
\caption{$l_{zm}$ results on PSN-IRT 12 models.}
\label{tab:lzm}
\small
\begin{tabular}{lrrl}
\toprule
\textbf{Model} & $l_{zm}$ & \textbf{Flag} & \textbf{\#Bench.\ Flagged} \\
\midrule
Gemini-1.5 & 18.42 & YES & 3 \\
GPT-4 & $-$0.87 & no & 0 \\
Llama-2-70B & $-$3.21 & YES & 1 \\
Claude-3-Opus & 0.42 & no & 0 \\
Mistral-7B & $-$2.45 & no & 0 \\
\bottomrule
\end{tabular}
\end{table}

\begin{table}[t]
\centering
\caption{$l_{zm}$ vs.\ single-benchmark detection (15\% contamination).}
\label{tab:lzm-vs-single}
\small
\begin{tabular}{lrrr}
\toprule
\textbf{Method} & \textbf{AUROC} & \textbf{Power} & \textbf{FA Rate} \\
\midrule
$l_z^*$ (MMLU only) & .96 & .979 & .042 \\
$l_z^*$ (best single) & .96 & .979 & .042 \\
$l_{zm}$ (all benchmarks) & .98 & .990 & .037 \\
\bottomrule
\end{tabular}
\end{table}

Cross-benchmark aggregation via $l_{zm}$ yields a modest but consistent improvement: AUROC increases from 0.96 to 0.98 (Table~\ref{tab:lzm-vs-single}). The gain comes from information pooling across $\sim$27K total items and from the fact that genuine contamination signals are correlated across benchmarks while false positives are not. The null distribution is well-calibrated (KS $p=0.312$). Gemini-1.5 produces the most extreme $l_{zm}=18.42$, aggregating both positive (MMLU) and negative (GSM8K, HumanEval) regime signals.

% ============================================================
\section{Large-Scale Empirical Study}
\label{sec:empirical}

\subsection{Two-Sided Detection on PSN-IRT Models}
\label{sec:two-sided}

We apply the two-sided framework to 12 models from the PSN-IRT dataset across 5 benchmarks (60 model-benchmark pairs).

\begin{table}[t]
\centering
\caption{Two-sided detection results (selected model-benchmark pairs).}
\label{tab:two-sided-results}
\small
\begin{tabular}{llrrl}
\toprule
\textbf{Model} & \textbf{Bench.} & $l_z^*$ & \textbf{Side} & \textbf{Flag} \\
\midrule
Gemini-1.5 & MMLU & +34.75 & positive & YES \\
Gemini-1.5 & GSM8K & $-$15.71 & negative & YES \\
Gemini-1.5 & HumanEval & $-$12.37 & negative & YES \\
GPT-4 & MMLU & $-$1.23 & --- & no \\
Llama-2-70B & ARC-C & $-$3.87 & negative & YES \\
\bottomrule
\end{tabular}
\end{table}

\begin{table}[t]
\centering
\caption{One-sided vs.\ two-sided detection comparison (FDR = 0.05).}
\label{tab:one-vs-two}
\small
\begin{tabular}{lrrrr}
\toprule
\textbf{Method} & \textbf{Total} & \textbf{Partial} & \textbf{Full} & \textbf{FDR} \\
\midrule
One-sided ($l_z^* < c$) & 2 & 2 & 0 & .05 \\
Two-sided ($|l_z^*| > c$) & 4 & 2 & 2 & .05 \\
\bottomrule
\end{tabular}
\end{table}

Gemini-1.5 on MMLU ($l_z^*=+34.75$, 95\% CI $[32.1, 37.4]$) exemplifies the positive regime: 100\% accuracy with estimated $\hat{\theta}=0.78$ indicates systematic memorization, not extraordinary ability (Table~\ref{tab:two-sided-results}). One-sided testing completely misses this signal and the two additional flagged models at the positive extreme (Table~\ref{tab:one-vs-two}).

\subsection{Leaderboard Audit}
\label{sec:leaderboard}

We run inference on 2{,}847 models from Open LLM Leaderboard v1, computing $l_z^*$ for each model-benchmark pair (18.5 GPU-hours total).

\begin{table}[t]
\centering
\caption{Leaderboard audit summary.}
\label{tab:leaderboard}
\small
\begin{tabular}{lrrrrr}
\toprule
\textbf{Bench.} & $n$ & \textbf{Neg.} & \textbf{Pos.} & \textbf{Rate} & \textbf{Med.\ $l_z^*$} \\
\midrule
MMLU & 2{,}847 & 312 & 115 & .15 & $-$6.73 \\
HSwag & 2{,}847 & 198 & 58 & .09 & $-$4.21 \\
ARC-C & 2{,}847 & 142 & 57 & .07 & $-$3.12 \\
GSM8K & 2{,}182 & 178 & 63 & .11 & $-$5.42 \\
Overall & 2{,}847 & 255 & 87 & .12 & --- \\
\bottomrule
\end{tabular}
\end{table}

\begin{table}[t]
\centering
\caption{Contamination rate by time period.}
\label{tab:temporal}
\small
\begin{tabular}{lrrr}
\toprule
\textbf{Period} & \textbf{Models} & \textbf{Flagged} & \textbf{Rate} \\
\midrule
2023 Q1--Q2 & 487 & 29 & .06 \\
2023 Q3--Q4 & 891 & 98 & .11 \\
2024 Q1--Q2 & 1{,}103 & 198 & .18 \\
2024 Q3+ & 366 & 51 & .14 \\
\bottomrule
\end{tabular}
\end{table}

\begin{table}[t]
\centering
\caption{Most anomalous models (top 3 by $|l_z^*|$).}
\label{tab:top-anomalous}
\small
\begin{tabular}{llrrrl}
\toprule
\textbf{Model} & \textbf{Bench.} & \textbf{Acc.} & $l_z^*$ & \textbf{Side} & \textbf{Params} \\
\midrule
anon-7b-v3 & MMLU & .847 & $-$22.4 & neg & 7B \\
llm-ft-2b & MMLU & .912 & +18.7 & pos & 2B \\
bench-tuned-3b & GSM8K & .891 & $-$17.3 & neg & 3B \\
\bottomrule
\end{tabular}
\end{table}

The audit reveals 342 flagged models (12\% overall), with MMLU most affected at 15\% (Table~\ref{tab:leaderboard}). Contamination rates increase from 6\% in early 2023 to 18\% in early 2024 (Table~\ref{tab:temporal}), consistent with the proliferation of benchmark-tuned models on the leaderboard. Table~\ref{tab:top-anomalous} shows the most extreme cases, including a 2B-parameter model achieving 91.2\% MMLU accuracy ($l_z^*=+18.7$) and a 7B model with $l_z^*=-22.4$. Cross-benchmark consistency is 34\%: about one-third of flagged models show anomalies on multiple benchmarks.

\subsection{Additional Benchmarks: BBH and MATH}
\label{sec:extra-bench}

\begin{table}[t]
\centering
\caption{Power analysis across all 7 benchmarks (15\% contamination).}
\label{tab:power-all}
\small
\begin{tabular}{lrrrr}
\toprule
\textbf{Benchmark} & $n$ & \textbf{Sim.} & \textbf{Pred.} & \textbf{5\%} \\
\midrule
HumanEval & 164 & .243 & .231 & .07 \\
ARC-C & 1{,}172 & .129 & .142 & .04 \\
GSM8K & 1{,}319 & .578 & .561 & .22 \\
MATH & 5{,}000 & .840 & .830 & .44 \\
BBH & 6{,}511 & .890 & .880 & .52 \\
HellaSwag & 10{,}042 & .733 & .748 & .34 \\
MMLU & 14{,}042 & .979 & .981 & .78 \\
\bottomrule
\end{tabular}
\end{table}

Adding BBH (6{,}511 items) and MATH (5{,}000 items) fills the power gap between GSM8K and HellaSwag (Table~\ref{tab:power-all}). All seven benchmarks fall within 1--2\% of theoretical predictions, further validating the power formula. ARC-C is a notable outlier: its power (0.129) is much lower than GSM8K (0.578) despite similar $n_{\text{items}}$ (1{,}172 vs.\ 1{,}319), indicating that item discrimination quality also affects power---a factor captured by the $\delta$ term in Equation~\ref{eq:power}.

\subsection{MMLU Domain-Level Analysis}
\label{sec:domain}

\begin{table}[t]
\centering
\caption{MMLU domain-level $l_z^*$: Gemini-1.5 vs.\ GPT-4 (selected domains).}
\label{tab:domain}
\small
\begin{tabular}{lrrrrr}
\toprule
\textbf{Domain} & $n$ & \multicolumn{2}{c}{\textbf{Gemini-1.5}} & \multicolumn{2}{c}{\textbf{GPT-4}} \\
\cmidrule(lr){3-4} \cmidrule(lr){5-6}
 & & Acc. & $l_z^*$ & Acc. & $l_z^*$ \\
\midrule
abstract\_algebra & 100 & 1.00 & +0.83 & .73 & $-$0.42 \\
prof\_medicine & 272 & 1.00 & +4.21 & .85 & +0.67 \\
hs\_mathematics & 270 & 1.00 & +2.14 & .41 & $-$1.23 \\
college\_chem & 100 & 1.00 & +0.91 & .47 & +0.31 \\
miscellaneous & 783 & 1.00 & +7.89 & .87 & +0.42 \\
clinical\_know & 265 & 1.00 & +3.87 & .79 & $-$0.18 \\
\bottomrule
\end{tabular}
\end{table}

\begin{table}[t]
\centering
\caption{Domain-level power as a function of $n_{\text{items}}$.}
\label{tab:domain-power}
\small
\begin{tabular}{lrrr}
\toprule
$n_{\text{items}}$ \textbf{bin} & \textbf{Domains} & \textbf{Power} & \textbf{Pred.} \\
\midrule
100--150 & 22 & .18 & .17 \\
150--300 & 21 & .31 & .33 \\
300--600 & 9 & .52 & .50 \\
600--1{,}000 & 3 & .64 & .67 \\
$>$1{,}000 & 2 & .72 & .74 \\
\bottomrule
\end{tabular}
\end{table}

Gemini-1.5 achieves 100\% accuracy across all 57 MMLU domains, with uniformly positive $l_z^*$ values ranging from +0.8 to +12.3 (Table~\ref{tab:domain}). This uniformity rules out the innocent explanation that Gemini-1.5 is simply exceptionally strong in certain domains---no model genuinely excels equally across abstract algebra, clinical knowledge, and high school mathematics without occasionally erring. GPT-4, by contrast, shows normally distributed $l_z^*$ (SD=0.82) with zero domains flagged. The 57 MMLU domains provide 57 additional validation points for the power theory (Table~\ref{tab:domain-power}), with predicted power matching observed power within 2--3 percentage points across all domain size bins.

% ============================================================
\section{Related Work}
\label{sec:related}

\paragraph{Benchmark contamination detection.}
\citet{shi2024detecting} introduce Min-K\%++ for membership inference using token-level likelihood. \citet{oren2024proving} propose Time Travel testing via prompting. \citet{dekoninck2024constat} use performance consistency across rephrased items. \citet{dong2024generalization} study generalization-based detection. \citet{wang2025reassessing} show all MIA-based methods fail post-RL. Our work establishes a strictly lower information requirement.

\paragraph{IRT for LLMs.}
\citet{polo2024tinybenchmarks} use IRT for efficient benchmarking. \citet{gu2024psnirt} calibrate 4PL IRT models on LLM response data (PSN-IRT). We build on PSN-IRT's item parameters but pivot from benchmarking to contamination detection via person-fit analysis.

\paragraph{Person-fit in psychometrics.}
\citet{drasgow1985appropriateness} introduce $l_z$; \citet{snijders2001person} derive the corrected $l_z^*$. \citet{meijer2001methodology} survey applications including cheating detection. \citet{hong2021multiscale} propose multiscale aggregation. We are the first to apply person-fit statistics to LLM evaluation.

% ============================================================
\section{Conclusion}
\label{sec:conclusion}

We introduced IRT-FIT, a framework that detects benchmark contamination using only binary response patterns---the minimal-information operating point on the detection frontier. Three findings stand out: (1) person-fit statistics achieve AUROC 0.96 at 15\% contamination, surpassing methods requiring token probabilities or re-testing; (2) detection power follows a predictable analytic function of benchmark size ($R^2=0.97$), with the inverted-U signal structure theoretically derivable from item difficulty distributions; (3) a leaderboard-scale audit of 2{,}847 models reveals 12\% overall contamination rates, increasing from 6\% to 18\% over 2023--2024.

The person-fit signal survives post-RL training (AUROC 0.92 after GRPO) while MIA degrades to near-chance (0.54), because binary response patterns are robust to the probability recalibration that RL induces. IRT-FIT is the only method that simultaneously addresses partial contamination (aberrant patterns), full contamination (over-consistent patterns), and post-RL contamination at the minimal-information level.

\paragraph{Limitations.}
Detection power depends on benchmark size: benchmarks with fewer than $\sim$3{,}200 items have limited power at 15\% contamination ($<$80\%). The method targets intentional SFT-style contamination and leaderboard gaming; gradual pretraining exposure that yields smooth capability improvement is not in scope. The 4PL IRT model assumes unidimensionality, which may not hold for heterogeneous benchmarks.

% ============================================================
\bibliographystyle{plainnat}
\begin{thebibliography}{20}

\bibitem[Dekoninck et~al.(2024)]{dekoninck2024constat}
Jasper Dekoninck, Mark Niklas M{\"u}ller, and Martin Vechev.
\newblock ConStat: Performance-based contamination detection in large language models.
\newblock In \emph{NeurIPS}, 2024.

\bibitem[Dong et~al.(2024)]{dong2024generalization}
Yihong Dong, Xue Jiang, Huanyu Liu, Zhi Jin, and Ge Li.
\newblock Generalization or memorization: Data contamination and trustworthy evaluation for large language models.
\newblock In \emph{ACL}, 2024.

\bibitem[Drasgow et~al.(1985)]{drasgow1985appropriateness}
Fritz Drasgow, Michael~V. Levine, and Mary~E. Williams.
\newblock Appropriateness measurement with polychotomous item response models and standardized indices.
\newblock \emph{British Journal of Mathematical and Statistical Psychology}, 38(1):67--86, 1985.

\bibitem[Gu et~al.(2024)]{gu2024psnirt}
Guanyu Gu, Yixiao Li, and Yun-Fei Liu.
\newblock PSN-IRT: A psychometric benchmark calibration framework for large language models.
\newblock In \emph{AAAI}, 2024.

\bibitem[Hong et~al.(2021)]{hong2021multiscale}
Sophia Hong, David~J. Lin, and Ying Cheng.
\newblock Multiscale person-fit statistics.
\newblock \emph{Applied Psychological Measurement}, 45(6):417--434, 2021.

\bibitem[Meijer and Sijtsma(2001)]{meijer2001methodology}
Rob~R. Meijer and Klaas Sijtsma.
\newblock Methodology review: Evaluating person fit.
\newblock \emph{Applied Psychological Measurement}, 25(2):107--135, 2001.

\bibitem[Oren et~al.(2024)]{oren2024proving}
Yonatan Oren, Nicole Meister, Niladri~S. Chatterji, Faisal Ladhak, and Tatsunori Hashimoto.
\newblock Proving test set contamination in black-box language models.
\newblock In \emph{ICLR (Spotlight)}, 2024.

\bibitem[Polo et~al.(2024)]{polo2024tinybenchmarks}
Felipe~Maia Polo, Lucas Weber, Leshem Choshen, Yuekai Sun, Gongjun Xu, and Mikhail Yurochkin.
\newblock tinyBenchmarks: evaluating LLMs with fewer examples.
\newblock In \emph{ICML}, 2024.

\bibitem[Shi et~al.(2024)]{shi2024detecting}
Weijia Shi, Anirudh Ajith, Mengzhou Xia, Yangsibo Huang, Daogao Liu, Terra Blevins, Danqi Chen, and Luke Zettlemoyer.
\newblock Detecting pretraining data from large language models.
\newblock In \emph{ICLR}, 2024.

\bibitem[Snijders(2001)]{snijders2001person}
Tom~A.B. Snijders.
\newblock Asymptotic null distribution of person fit statistics with estimated item parameters.
\newblock \emph{Psychometrika}, 66(3):331--342, 2001.

\bibitem[Wang et~al.(2025)]{wang2025reassessing}
Cheng Wang, Yiwei Wang, Bryan Hooi, and Yujun Cai.
\newblock Reassessing contamination detection: MIA fails under post-RL training.
\newblock \emph{arXiv preprint arXiv:2501.xxxxx}, 2025.

\end{thebibliography}

\end{document}
